% Standalone solution dossier for prop:a4-local-excess.
\documentclass[../main.tex]{subfiles}
\begin{document}

% === SOLUTION HEADER (proof provenance) =====================================
% ledger-node: prop:a4-local-excess
% refines: prop:a4-local-excess
% bounded_by: obs:restricted-not-finite
% evidence_run: none
% checked_by: agent
% author: /root/a4_a5
% reviewer: /root/review_a4_a5
% dossier_editor: /root/review_a4_a5
% review: research/reviews/2026-08-21-a4-a5-independent-agent-audit.md
% date: 2026-08-21
% ============================================================================

\section*{Solution: optimizer-centred excess-KL constants}

\begin{theorem}
Let $\pi\in\mathcal P_2(\R^d)$ and let $q_\eta\in\mathcal P_2(\R^d)$, $\eta\in\R^p$, be an
identifiable family, so all displayed means are finite. Let $\eta_\star$ be the unique interior
minimizer of $K(\eta)=\KL(q_\eta\|\pi)$, put $\delta=K(\eta_\star)$ and
$q_\star=q_{\eta_\star}$. For some $\rho_0>0$, assume the complete sublevel
$\{\eta:K(\eta)\le\delta+\rho_0\}$ is contained in one compact parameter chart on which $K$
is continuous, and assume the minimizer is isolated there: every neighbourhood $V$ of
$\eta_\star$ satisfies
\[
 \inf_{\eta\notin V,\ K(\eta)\le\delta+\rho_0}K(\eta)>\delta.
\]
Suppose also that the following expansions hold uniformly as $h\to0$:
\[
 2\{K(\eta_\star+h)-\delta\}=h^\top F_\star h+O(\norm h^3),\qquad F_\star\succ0,
\]
\[
 W_2^2(q_{\eta_\star+h},q_\star)=h^\top G_\star h+O(\norm h^3),
 \qquad \E_{q_{\eta_\star+h}}\theta-\E_{q_\star}\theta=B_\star h+O(\norm h^2).
\]
Then, as $\rho\downarrow0$,
\[
 \sup_{0<K(\eta)-\delta\le\rho}
 \frac{W_2^2(q_\eta,q_\star)}{2\{K(\eta)-\delta\}}
 =\lambda_{\max}(F_\star^{-1/2}G_\star F_\star^{-1/2})+O(\sqrt\rho),
\]
with the analogous mean constant equal to
$\lambda_{\max}(F_\star^{-1/2}B_\star^\top B_\star F_\star^{-1/2})+O(\sqrt\rho)$.
\end{theorem}

\begin{proof}
The compact isolation and KL expansion imply $\norm h=O(\sqrt\rho)$ on the excess sublevel.
After subtracting $\delta$, both numerator and denominator vanish quadratically. Their ratio is
the corresponding generalized Rayleigh quotient plus a uniform $O(\norm h)=O(\sqrt\rho)$
error. Maximization gives the upper bound, and arbitrarily small multiples of a top generalized
eigenvector give the lower bound. The squared mean expansion is
$h^\top B_\star^\top B_\star h+O(\norm h^3)$ and is handled identically.
\end{proof}

\paragraph{Scope.}
This controls optimization error relative to $q_\star$. It does not remove the approximation
error $W_2(q_\star,\pi)$.

\end{document}
